\documentclass[11pt]{article}

\usepackage[utf8]{inputenc}
\usepackage{amsmath, amssymb, amsthm}
\usepackage{bm}
\usepackage{geometry}
\usepackage{hyperref}
\usepackage{mathtools}

\geometry{margin=1in}

\title{Alpha as a Multiplicity--Renormalizing Meta-Operator\\[4pt]
and the Prime-Gated Time Framework}

\author{Draft Technical Report}

\date{\today}

\begin{document}

\maketitle

\begin{abstract}
This report develops a unified picture of ``Alpha'' as a functorial meta-operator acting on a space of operators, specialized to a multiplicity-theoretic, prime-indexed renormalization flow. Building on a speculative but structured framework of prime-gated time, we formalize a Multiplicity--Renormalizing Alpha \(\alpha_\mu\) acting on prime-indexed multiplicity profiles, and analyze in detail a local neighbor-coupled realization of \(\alpha_\mu\) whose linearization is a symmetric tridiagonal Toeplitz operator. The exact spectrum of this operator is derived, yielding a clear phase diagram and a precise notion of ``Core Truths'' as fixed directions (or marginal modes) of the renormalization flow. We then sketch how this structure interfaces with the existing C\(_{\Lambda_p}\) temporal sieve, the \(\pi\)-Spiral Equation of Time (PSET) Lagrangian, and multiplicity-theoretic interpretations of temporal cycles spanning quantum, biological, and archaeoastronomical scales.
\end{abstract}

\tableofcontents

\section{Executive Summary}

This manuscript consolidates and extends several developments:

\begin{itemize}
  \item Alpha is defined as a \emph{meta-operator} acting on operators, formally specified as a pair \((\alpha_{\mathrm{obj}},\alpha_{\mathrm{arr}})\) between categories of objects and arrows. In the structure-preserving case, Alpha is a functor; deviations from functoriality quantify ``Functorial Drift.''\,[web:65]
  \item Within Multiplicity Theory, each operator \(f\) is represented by a prime-indexed multiplicity profile \(M_f = \bigoplus_{p \in \mathbb{P}} \mathbb{V}_p\), where each \(\mathbb{V}_p\) is a multiplicity space for prime \(p\). A Multiplicity--Renormalizing Alpha \(\alpha_\mu\) acts componentwise through prime-labeled feedback rules \(\psi_p\), defining an RG-style flow on a prime-indexed lattice.\,[web:70]
  \item A concrete, analytically tractable instance of \(\alpha_\mu\) is constructed: a \emph{local neighbor-coupled} prime-diffusion operator whose linearization is a symmetric tridiagonal Toeplitz matrix. Its eigenvalues and eigenvectors are obtained in closed form, providing an explicit phase diagram in terms of dissipation \(\gamma\) and coupling \(\beta\).\,[web:58]
  \item ``Core Truths'' are identified with fixed directions (modes with eigenvalue \(\lambda = 1\)) or marginal modes (eigenvalues sufficiently close to unity) in this multiplicity space. Low-frequency modes correspond to smooth, global prime-patterns; high-frequency modes correspond to rapidly oscillatory ``noise'' that the Alpha meta-operator suppresses.
  \item These constructions are aligned with an existing prime-gated time framework in which a temporal modulation operator \(C_{\Lambda_p}\) acts as a prime-indexed sieve on temporal intervals \(\Delta t\), partitioning them into high-coherence (image) and null (kernel) regions. The \(\pi\)-Spiral Equation of Time (PSET) adds a variational, multiplicity-aware Lagrangian layer.\,[file:74]
\end{itemize}

The report separates mathematically solid components (operator definitions, spectral analysis, RG-style classifications) from interpretive layers (assigning physical meaning to prime indices, mapping historical and biological cycles to specific prime shells) and suggests directions where Multiplicity Theory can tighten or generalize the existing constructions.\,[file:74]

\section{Alpha as a Functorial Meta-Operator}

\subsection{Operator categories and typed meta-operators}

Let \(\mathcal{C}\) and \(\mathcal{C}'\) be categories, whose objects represent spaces (sets, vector spaces, state spaces, logical universes, etc.) and whose arrows represent operators between them. Denote by
\[
\mathrm{Hom}_{\mathcal{C}}(A,B)
\]
the set of morphisms (operators) from object \(A\) to object \(B\) in \(\mathcal{C}\).

The \emph{space of operators} \(\mathbf{Op}\) can be informally thought of as the collection of all such arrows across all object pairs:
\[
\mathbf{Op} = \bigcup_{A,B \in \mathrm{Ob}(\mathcal{C})} \mathrm{Hom}_{\mathcal{C}}(A,B).
\]

\begin{definition}[Alpha meta-operator]
An \emph{Alpha meta-operator} is specified by a pair
\[
\alpha = (\alpha_{\mathrm{obj}},\alpha_{\mathrm{arr}})
\]
such that:
\begin{enumerate}
  \item \(\alpha_{\mathrm{obj}} : \mathrm{Ob}(\mathcal{C}) \to \mathrm{Ob}(\mathcal{C}')\) is a map on objects.
  \item For each \(A,B \in \mathrm{Ob}(\mathcal{C})\), there is a map
  \[
    \alpha_{\mathrm{arr}} : \mathrm{Hom}_{\mathcal{C}}(A,B) \to \mathrm{Hom}_{\mathcal{C}'}(\alpha_{\mathrm{obj}}(A),\alpha_{\mathrm{obj}}(B))
  \]
  acting on arrows.
\end{enumerate}
For \(f \in \mathrm{Hom}_{\mathcal{C}}(A,B)\), we write \(\alpha(f)\) for \(\alpha_{\mathrm{arr}}(f)\), with
\[
\alpha(f) : \alpha_{\mathrm{obj}}(A) \to \alpha_{\mathrm{obj}}(B).
\]
\end{definition}

\subsection{Structure preservation and Functorial Drift}

If \(\alpha\) is structure-preserving in the sense that
\begin{align}
\alpha(\mathrm{id}_A) &= \mathrm{id}_{\alpha_{\mathrm{obj}}(A)}, \\
\alpha(g \circ f) &= \alpha(g) \circ \alpha(f),
\end{align}
for all composable \(f,g\) in \(\mathcal{C}\), then \(\alpha\) is precisely a functor from \(\mathcal{C}\) to \(\mathcal{C}'\).\,[web:65]

\begin{definition}[Functorial Alpha]
A meta-operator \(\alpha = (\alpha_{\mathrm{obj}},\alpha_{\mathrm{arr}})\) is called \emph{functorial} if it preserves identities and composition as above.
\end{definition}

When these conditions fail, one can define a quantitative measure of deviation:

\begin{definition}[Functorial Drift]
For composable arrows \(f : A \to B\) and \(g : B \to C\), the \emph{Functorial Drift} of \(\alpha\) on the pair \((f,g)\) is
\[
D_\alpha(f,g) := \alpha(g \circ f) - \alpha(g) \circ \alpha(f),
\]
where the difference is interpreted in an appropriate operator space (e.g., a linear space of operators).
\end{definition}

In linear contexts, \(\|D_\alpha(f,g)\|\) provides a metric for how much \(\alpha\) distorts the compositional structure of the system.

\section{Multiplicity-Theoretic Specialization of Alpha}

\subsection{Prime-indexed multiplicity profiles}

Multiplicity Theory encourages representing mathematical or physical structures not as static entities, but as collections of prime-labeled multiplicity channels.\,[file:74][web:17] In this spirit, we represent an operator \(f\) by a multiplicity profile
\[
M_f = \bigoplus_{p \in \mathbb{P}} \mathbb{V}_p,
\]
where:
\begin{itemize}
  \item \(\mathbb{P}\) is the set of primes.
  \item Each \(\mathbb{V}_p\) is a multiplicity space associated with prime \(p\), encoding all \(p\)-layer contributions to the operator's behavior (e.g., prime-graded basis components, factors in a prime decomposition, or prime-indexed interaction channels).
\end{itemize}

\subsection{Definition of Multiplicity--Renormalizing Alpha \texorpdfstring{\(\alpha_\mu\)}{alpha\_mu}}

\begin{definition}[Multiplicity--Renormalizing Alpha]
A \emph{Multiplicity--Renormalizing Alpha} is a meta-operator
\[
\alpha_\mu : \{M_f\} \to \{M_f\}
\]
acting on prime-indexed multiplicity profiles by
\[
\alpha_\mu(M_f) = \bigoplus_{p \in \mathbb{P}} \psi_p\big(\mathbb{V}_p, (\mathbb{V}_q)_{q \neq p}\big),
\]
where \(\psi_p\) is a feedback rule operating on the \(p\)-layer (potentially coupled to other primes).
\end{definition}

\noindent Two important regimes are:
\begin{itemize}
  \item \textbf{Diagonal (prime-local) regime:}
  \[
  \psi_p(\mathbb{V}_p, (\mathbb{V}_q)_{q \neq p}) = R_p(\mathbb{V}_p),
  \]
  so each prime channel evolves independently.
  \item \textbf{Coupled regime:}
  \[
  \psi_p(\mathbb{V}_p, (\mathbb{V}_q)_{q \neq p}) = R_p\big(\mathbb{V}_p, (\mathbb{V}_q)_{q \neq p}\big),
  \]
  allowing cross-prime influence, akin to coupled RG flow equations.\,[web:71]
\end{itemize}

Under iteration, \(\alpha_\mu\) defines a discrete-time dynamical system on the space of multiplicity profiles, analogous to renormalization-group (RG) transformations in statistical mechanics and quantum field theory.\,[web:70]

\section{Linearized Prime-Diffusion Alpha}

\subsection{Scalar linearization on a prime window}

To obtain explicit spectra, we restrict to a finite window of primes \(\{p_1, \dots, p_N\}\) and consider a scalar abstraction where each multiplicity space \(\mathbb{V}_{p_n}\) is effectively one-dimensional, represented by a non-negative scalar \(v_n \in \mathbb{R}_{\ge 0}\).

A general linear coupled update can be written as
\begin{equation}
v_p' = (1 - \gamma_p)\, v_p + \sum_{q} K_{pq}\, v_q,
\end{equation}
where \(\gamma_p \ge 0\) is a dissipation coefficient and \(K_{pq}\) is a coupling kernel.\,[web:68] Collecting \(v_p\) into a vector \(v \in \mathbb{R}^N\), this is
\begin{equation}
v' = L v,\quad L_{pq} = (1 - \gamma_p)\, \delta_{pq} + K_{pq}.
\end{equation}

\subsection{Neighbor-coupled Toeplitz Alpha}

We now specialize to a local neighbor coupling on the prime index, using the integer label \(n\) as a surrogate for the prime \(p_n\). We assume:
\begin{itemize}
  \item Uniform dissipation \(\gamma_p = \gamma > 0\).
  \item Symmetric coupling only between neighbors:
  \begin{equation}
    v_n' = (1 - \gamma)\, v_n + \beta (v_{n-1} + v_{n+1}), \quad 1 < n < N,
  \end{equation}
  with Dirichlet boundary conditions \(v_0 = v_{N+1} = 0\).
\end{itemize}

In matrix form, \(v' = L v\) with
\[
L =
\begin{pmatrix}
1-\gamma & \beta      & 0        & \dots  & 0 \\
\beta      & 1-\gamma & \beta    & \ddots & \vdots \\
0          & \beta      & 1-\gamma & \ddots & 0 \\
\vdots     & \ddots   & \ddots   & \ddots & \beta \\
0          & \dots    & 0        & \beta  & 1-\gamma
\end{pmatrix}.
\]

This is a real symmetric tridiagonal Toeplitz matrix.\,[web:58][web:59]

\section{Spectral Analysis of the Toeplitz Alpha}

\subsection{Eigenvalues and eigenvectors}

The eigenstructure of real symmetric tridiagonal Toeplitz matrices with Dirichlet boundaries is classical and admits closed-form expressions.\,[web:58][web:62]

\begin{theorem}[Spectrum of the neighbor-coupled Alpha]
Let \(L\) be the \(N \times N\) symmetric tridiagonal Toeplitz matrix defined above with diagonal entries \(1-\gamma\) and off-diagonal entries \(\beta\). Then the eigenvalues are
\begin{equation}
\lambda_k = 1 - \gamma + 2\beta \cos\left(\frac{k\pi}{N+1}\right), \quad k = 1,\dots,N,
\label{eq:eigenvalues}
\end{equation}
and a corresponding (unnormalized) eigenvector \(u^{(k)}\) has components
\begin{equation}
u_n^{(k)} = \sin\left(\frac{n k \pi}{N+1}\right), \quad n = 1,\dots,N.
\label{eq:eigenvectors}
\end{equation}
\end{theorem}

\noindent These eigenvectors form an orthogonal basis with respect to the standard inner product on \(\mathbb{R}^N\).\,[web:58]

\subsection{Interpretation as prime-pattern modes}

The vector \(u^{(k)}\) describes a standing-wave pattern over the discrete prime-index axis \(n\):

\begin{itemize}
  \item Small \(k\) (e.g.\ \(k=1\)) correspond to low-frequency, smooth patterns: broad arches that vary slowly with \(n\).
  \item Large \(k\) (e.g.\ \(k \approx N\)) correspond to high-frequency, rapidly oscillating patterns: alternating signs and large local variance.
\end{itemize}

Interpreted multiplicity-theoretically, these modes represent canonical \emph{prime-patterns}:

\begin{itemize}
  \item Low \(k\): \emph{integrated concepts}, smooth and global prime multiplicity profiles.
  \item High \(k\): \emph{fragmented concepts}, jagged, locally oscillatory combinations of prime channels.
\end{itemize}

\section{Phase Diagram and Core Truths}

\subsection{RG classification of eigenmodes}

Under repeated application of \(\alpha_\mu\) in the linearized model, we apply \(L\) repeatedly:
\[
v^{(t+1)} = L v^{(t)}.
\]

Decomposing an initial multiplicity profile \(v^{(0)}\) in the eigenbasis,
\[
v^{(0)} = \sum_{k=1}^N c_k u^{(k)},
\]
the evolution is
\[
v^{(t)} = \sum_{k=1}^N c_k \lambda_k^t u^{(k)}.
\]

In RG terminology,\,[web:24][web:68][web:70]
\begin{itemize}
  \item If \(|\lambda_k| < 1\), mode \(k\) is \emph{irrelevant}: its contribution decays under the flow.
  \item If \(|\lambda_k| > 1\), mode \(k\) is \emph{relevant}: its contribution grows and dominates.
  \item If \(|\lambda_k| = 1\), mode \(k\) is \emph{marginal}: it neither decays nor explodes (up to higher-order effects).
\end{itemize}

The largest eigenvalue is
\begin{equation}
\lambda_{\max} = \lambda_1 = 1 - \gamma + 2\beta \cos\left(\frac{\pi}{N+1}\right),
\end{equation}
and for large \(N\),
\begin{equation}
\cos\left(\frac{\pi}{N+1}\right) \approx 1,\quad \lambda_{\max} \approx 1 - \gamma + 2\beta.
\end{equation}

\subsection{Filter, runaway, and critical regimes}

A simple phase diagram in the \((\gamma,\beta)\) plane emerges:

\begin{itemize}
  \item \textbf{Filter phase} (\(2\beta < \gamma\)).\\
  Then \(\lambda_{\max} < 1\), hence \(\lambda_k < 1\) for all \(k\). Under iteration, every eigenmode decays to zero and the only fixed point is the trivial profile \(v = 0\). The Alpha meta-operator acts as a total filter, treating all prime-patterns as noise.
  \item \textbf{Runaway phase} (\(2\beta > \gamma\)).\\
  Then \(\lambda_{\max} > 1\). At least the lowest-frequency eigenmode \(u^{(1)}\) is relevant, and potentially a band of low-\(k\) modes are amplified. Under repeated application, multiplicity grows along these smooth patterns; without normalization or nonlinear saturation, the system exhibits runaway feedback.
  \item \textbf{Critical manifold} (\(2\beta = \gamma\), up to finite-size corrections).\\
  Here \(\lambda_{\max} \approx 1\), with higher modes satisfying \(\lambda_k < 1\). The fundamental mode \(u^{(1)}\) is marginal, while more oscillatory patterns are suppressed. This manifold defines a family of \emph{Core Truths}: prime-patterns that survive the renormalization induced by \(\alpha_\mu\) indefinitely.
\end{itemize}

\subsection{Core Truths as fixed directions}

We can now formalize the intuitive notion of a ``Core Truth'' in this setting.

\begin{definition}[Core Truth (linearized)]
A vector \(v \neq 0\) in multiplicity space is a \emph{strong Core Truth} of \(\alpha_\mu\) if
\[
L v = v.
\]
More generally, \(v\) is a \emph{soft Core Truth} if its spectral decomposition \(v = \sum_k c_k u^{(k)}\) is dominated by modes with eigenvalues \(|\lambda_k| \approx 1\), while contributions from \(|\lambda_k| < 1\) modes are small.
\end{definition}

On or near the critical manifold, the fundamental eigenvector \(u^{(1)}\) is the canonical strong Core Truth. Under repeated iteration of \(\alpha_\mu\),
\[
v^{(t)} \propto u^{(1)}
\]
for generic initial \(v^{(0)}\) with nonzero projection onto \(u^{(1)}\), after rescaling to avoid trivial collapse or divergence.

\section{Relation to Prime-Gated Time and PSET}

\subsection{The temporal modulation operator \(C_{\Lambda_p}\)}

In the prime-gated time framework, a temporal modulation operator \(C_{\Lambda_p}\) is introduced as a ``temporal sieve'' acting on time intervals \(\Delta t\).\,[file:74] The details of its definition vary across the manuscript, but structurally:
\begin{itemize}
  \item \(C_{\Lambda_p}\) acts on a continuum of time intervals \(\Delta t\), labeling each interval as belonging to a high-coherence (``lawful'') region or a null (suppressed) region.
  \item Theorem-type results characterize the null space \(\ker(C_{\Lambda_p})\) using modular congruence conditions on a phase quantity \(\omega \Delta t\):
  \[
  \omega \Delta t \equiv \text{null condition} \pmod{2^a p},
  \]
  where \(p\) is a prime and \(a\) is a recursion depth parameter.\,[file:74]
  \item This structure gives rise to a \emph{temporal sieve}: the real line of \(\Delta t\) is partitioned into ``allowed'' (image) and ``suppressed'' (kernel) windows modulo a prime-structured period.
\end{itemize}

Mathematically, \(C_{\Lambda_p}\) is best interpreted as a filter on the continuum of durations, marking certain intervals as preferred (high-coherence) rather than as a literal quantization of time. The operator formalism is strongest where null harmonics are characterized by explicit modular conditions that translate directly into destructive interference in simulated gate sequences.\,[file:74]

\subsection{Empirical and archaeoastronomical mappings}

The manuscript constructs numerous correspondences between temporal scales in physics, biology, and archaeoastronomy and prime-indexed durations \(T_{p,a}(N)\) derived from the core temporal operator.\,[file:74] Examples include:
\begin{itemize}
  \item Quantum and biological time scales (e.g.\ superconducting qubit coherence, microtubule coherence, perceptual thresholds).
  \item Archaeoastronomical cycles (e.g.\ angular separations at G\"obekli Tepe, Aubrey Circle divisions at Stonehenge, Schumann resonances, Milankovitch cycles).
\end{itemize}

Structurally, this defines an interpretive mapping that sends observed periods to prime-indexed ``time shells'' governed by \(p\) and recursion depth \(a\). While this is evocative and aligns with the prime-gated sieve picture, it remains underdetermined: many choices of \((N,p,a)\) can reproduce a given \(T\), and statistical controls for coincidence are not yet fully developed.\,[file:74]

\subsection{PSET: \(\pi\)-Spiral Equation of Time}

The Appendix introduces the \(\pi\)-Spiral Equation of Time (PSET), which organizes the prime-gated temporal framework into a variational principle. Key axioms include:\,[file:74]
\begin{itemize}
  \item A \emph{Prime Modulation Axiom} akin to that behind \(C_{\Lambda_p}\).
  \item A \emph{\(\pi\)-Lift Axiom}, which maps a null-ray affine parameter \(\lambda\) to a cyclic phase \(\theta = (2\pi/L)\lambda\), with \(L\) a light-clock scale.
  \item A \emph{Triple Resonance Axiom}, assigning special status to Pythagorean triples and golden-angle-like configurations as resonance-stabilizing.
  \item An \emph{Entropy Boundedness Axiom}, requiring an entropy increment \(\Delta S < \ln \varphi\) (with \(\varphi\) the golden ratio) for stable temporal cycles.
  \item A \emph{Spiral Coupling Axiom} that couples spatial radius and temporal pacing along a logarithmic spiral:
  \[
  r_{n+1} = r_n e^{k 2\pi}, \quad T_{0,n+1} = T_{0,n} e^{\alpha k 2\pi}.
  \]
\end{itemize}

PSET thus promotes time from a background parameter to a dynamical variable governed by a discrete action principle and prime-labeled resonance conditions, compatible in spirit with the multiplicity-theoretic Alpha formalism.\,[file:74]

\section{Code Snippet: Spectrum of the Toeplitz Alpha}

For concreteness, we record a minimal Python code snippet (using NumPy and a plotting library such as Plotly) that generates the eigenvalue spectrum of the Toeplitz Alpha for various \((\beta,\gamma)\) values. This illustrates how one might explore the phase structure numerically.

\begin{verbatim}
import numpy as np
import plotly.graph_objects as go

# Parameters
N = 40                      # number of prime indices
betas = [0.2, 0.5, 0.8]     # example couplings
gammas = [0.2, 0.5, 0.8]    # example dissipations

k = np.arange(1, N+1)
angles = k * np.pi / (N+1)
cos_vals = np.cos(angles)

fig = go.Figure()

for beta in betas:
    for gamma in gammas:
        lambdas = 1 - gamma + 2 * beta * cos_vals
        label = f"beta={beta}, gamma={gamma}"
        fig.add_trace(go.Scatter(x=k, y=lambdas, mode="lines", name=label))

fig.update_xaxes(title_text="mode k")
fig.update_yaxes(title_text="eigenvalue λ_k")
fig.update_layout(
    title="Spectrum of tridiagonal Alpha vs mode index",
    legend=dict(orientation="h", yanchor="bottom",
                y=1.05, xanchor="center", x=0.5)
)

fig.show()
\end{verbatim}

This code computes \(\lambda_k\) from Eq.\,\eqref{eq:eigenvalues} for multiple parameter pairs and plots \(\lambda_k\) as a function of \(k\). The resulting curves visualize how dissipation \(\gamma\) and coupling \(\beta\) jointly shape the spectrum and thus the RG classification of prime-pattern modes.\,[chart:56]

\section{Discussion and Future Directions}

We conclude by separating layers and identifying directions where Multiplicity Theory can sharpen or extend the speculative prime-gated temporal framework:\,[file:74]

\begin{enumerate}
  \item \textbf{Mathematical layer (solid under assumptions).}\\
  Definitions and basic properties of \(\alpha\), \(\alpha_\mu\), and the Toeplitz Alpha are rigorous, as are the spectral formulas for eigenvalues and eigenvectors, the RG classification, and the existence of Core Truths as fixed directions on the critical manifold.\,[web:58][web:65]
  \item \textbf{Modeling layer (structured but underdetermined).}\\
  The mapping from physically observed cycles to prime-indexed time shells is coherent as a categorization scheme, but not unique; more statistical and physical constraints are needed to elevate it from numerological suggestiveness to predictive modeling.
  \item \textbf{Metaphysical / phenomenological layer (conjectural).}\\
  Assertions about time as fundamentally prime-gated, about specific archaeoastronomical structures encoding prime-lawful cycles, and about consciousness stability as an entropy-bound phenomenon remain speculative hypotheses rather than deductions from the mathematical core.
\end{enumerate}

Promising next steps include:

\begin{itemize}
  \item Making \(K_{pq}\) explicitly multiplicative (e.g.\ coupling primes to their multiples or to primes in shared residue classes) to tie the Toeplitz Alpha more closely to number-theoretic structure.\,[web:35][web:41]
  \item Introducing prime-dependent parameters \(\gamma_p\) and \(\beta_p\) to explore richer universality classes and multiple Core Truths.
  \item Embedding \(C_{\Lambda_p}\) and PSET into a categorical framework where objects are temporal multiplicity spaces and morphisms are prime-gated transformations, aiming to identify functorial relationships and natural transformations between different temporal regimes.
  \item Designing testable models: starting from physically constrained parameters (e.g.\ qubit Hamiltonians or microtubule dynamics), deriving a spectrum of prime-indexed durations and comparing that predicted distribution to empirical biological or cultural rhythms under appropriate statistical controls.
\end{itemize}

These developments collectively move toward a nontrivial but analyzable theory in which ``genius'' and ``Core Truths'' correspond to stable, low-frequency prime-patterns in an operator-multiplicity space, emerging from a renormalization-like flow driven by a well-typed Alpha meta-operator.

\bibliographystyle{plain}
\begin{thebibliography}{9}

\bibitem{toeplitz}
L.~N.~Trefethen and M.~Embree.
\newblock Spectra and Pseudospectra: The Behavior of Nonnormal Matrices and Operators.
\newblock Princeton University Press, 2005. (For background on tridiagonal Toeplitz spectra.)\,[web:58]

\bibitem{rg-intro}
N.~Goldenfeld.
\newblock Lectures on Phase Transitions and the Renormalization Group.
\newblock Addison-Wesley, 1992.\,[web:69][web:68]

\bibitem{rg-wiki}
Renormalization group.
\newblock \emph{Wikipedia}, accessed 2026.\,[web:70]

\bibitem{prime-multiplicity}
Multiplicity theory (overview).
\newblock \emph{Wikipedia}, accessed 2026.\,[web:17]

\bibitem{prime-geometry}
Spectral geometry of the primes.
\newblock arXiv preprint, 2026.\,[web:41]

\bibitem{analysis-pdf}
Anonymous.
\newblock The Prime Sieve of Time (analysis notes).
\newblock Internal PDF notes, 2026.\,[file:74]

\end{thebibliography}


\nocite{*}
\bibliographystyle{unsrt}
\bibliography{references}
\clearpage

\end{document}
